\paragraph{Definition and Derivation of Phase Kickback}\label{ex:phase-kickback-definition-and-derivation}

At this point we have already derived the mathematical result:
\begin{equation*}
    CU \ket{+}\ket{v} = \frac{\ket{0} + e^{i\phi}\ket{1}}{\sqrt{2}} \otimes \ket{v}
\end{equation*}
Now we can finally give a precise definition of \textbf{phase kickback}.

\highspace
\begin{definitionbox}[: Phase Kickback]
    \definition{Phase Kickback} is the quantum phenomenon in which the eigenvalue phase generated by applying a controlled unitary to one of its eigenstates appears as a relative phase on the control qubit instead of remaining on the target qubit.

    \highspace
    Equivalently, mathematically we can define phase kickback as the following property of controlled unitary operators:
    \begin{equation}\label{eq:phase-kickback-definition}
        CU\ket{+}\ket{v} = \frac{\ket{0} + e^{i\phi}\ket{1}}{\sqrt{2}} \otimes \ket{v}
    \end{equation}
    Where $U$ is a unitary operator, $\ket{v}$ is an eigenstate of $U$ and $e^{i\phi}$ is the corresponding eigenvalue of $U$, i.e.:
    \begin{equation*}
        U\ket{v} = e^{i\phi}\ket{v}
    \end{equation*}
    Notice the remarkable result:
    \begin{itemize}
        \item The target state is still \ket{v};
        \item The phase now modifies the control qubit $\ket{+}$ instead of the target qubit $\ket{v}$.
    \end{itemize}

    \highspace
    Formally, start with:
    \begin{equation*}
        \ket{\psi_{\text{in}}} = \ket{+} \ket{v}
    \end{equation*}
    Expanding the superposition of the control qubit \ket{+} we have:
    \begin{equation*}
        \ket{\psi_{\text{in}}} = \frac{\ket{0}\ket{v} + \ket{1}\ket{v}}{\sqrt{2}}
    \end{equation*}
    Now we can apply the controlled unitary $CU$ to the input state:
    \begin{align*}
        CU\ket{\psi_{\text{in}}} &= CU\left(\frac{\ket{0}\ket{v} + \ket{1}\ket{v}}{\sqrt{2}}\right) \\
        &= \frac{CU\ket{0}\ket{v} + CU\ket{1}\ket{v}}{\sqrt{2}} \\
        &= \frac{\ket{0}\ket{v} + \ket{1}U\ket{v}}{\sqrt{2}} \\
        &= \frac{\ket{0}\ket{v} + \ket{1} e^{i\phi}\ket{v}}{\sqrt{2}} \\
        &= \frac{\ket{0} + e^{i\phi}\ket{1}}{\sqrt{2}} \otimes \ket{v}
    \end{align*}
\end{definitionbox}

\highspace
\begin{flushleft}
    \textcolor{Green3}{\faIcon{question-circle} \textbf{Why is it called ``\emph{kickback}''?}}
\end{flushleft}
If we looked only at the target qubit \ket{v}, we would expect that the eigenvalue phase $e^{i\phi}$ would be applied to it:
\begin{equation*}
    U\ket{v} = e^{i\phi}\ket{v}
\end{equation*}
So one might think: ``\emph{the phase belongs to the target}''. However, after applying the controlled operation we obtain:
\begin{equation*}
    \left(\frac{\ket{0} + e^{i\phi}\ket{1}}{\sqrt{2}}\right) \otimes \ket{v}
\end{equation*}
The target has returned to exactly the same state \ket{v} as before, while the control now carries the phase. It therefore \textbf{looks as if} the \textbf{phase has travelled backwards from the target to the control}. \hl{This apparent transfer is why the phenomenon is called phase kickback.}

\highspace
\begin{flushleft}
    \textcolor{Red2}{\faIcon{exclamation-triangle} \textbf{The 4 conditions for phase kickback}}
\end{flushleft}
Phase kickback \textbf{does not happen automatically}. It requires four ingredients:
\begin{enumerate}
    \item \important{A controlled unitary}. The operation must be $CU$. If we simply apply $U$ to a qubit, there is no control branch where a relative phase can appear.
    
    
    \item \important{The target must be an eigenstate}. The target must satisfy:
    \begin{equation*}
        U\ket{v} = e^{i\phi}\ket{v}
    \end{equation*}
    Otherwise, $U$ changes the state itself, not just its phase. For example:
    \begin{equation*}
        X\ket{0} = \ket{1}
    \end{equation*}
    Is \textbf{not} an eigenstate relation. There is no simple phase to kick back because the target state changes from \ket{0} to \ket{1}.
    
    
    \item \important{The control must be in superposition}. The control must contain at least two branches, typically:
    \begin{equation*}
        \ket{+} = \frac{\ket{0} + \ket{1}}{\sqrt{2}}
    \end{equation*}
    Without superposition, there is only one branch, so every phase is global.
    
    
    \item \important{The eigenvalue must contribute a nontrivial phase}. If the eigenvalue is $1$, then the phase is trivial and there is no relative phase to kick back. For example, if $e^{i\phi} = 1$, then:
    \begin{equation*}
        \frac{\ket{0} + \ket{1}}{\sqrt{2}} = \frac{\ket{0} + 1\cdot\ket{1}}{\sqrt{2}}
    \end{equation*}
    Nothing changes. So the interesting cases are $e^{i\phi} \neq 1$.
\end{enumerate}
\textcolor{Green3}{\faIcon{question-circle} \textbf{Does the phase literally move?}} When we say ``\emph{the phase is kicked back}'', we do \textbf{not} mean that some physical object travels from one qubit to another. Nothing is moving, instead the mathematics of the tensor product allows us to rewrite:
\begin{equation*}
    \ket{0} \otimes \ket{v} + e^{i\phi}\ket{1} \otimes \ket{v}
\end{equation*}
As:
\begin{equation*}
    \left(\ket{0} + e^{i\phi}\ket{1}\right) \otimes \ket{v}
\end{equation*}
The phase was always a scalar multiplying one branch of the joint quantum state. So the ``kickback'' is a \textbf{mathematical redistribution of the phase in the tensor-product expression}, not the movement of a physical quantity.